\section{Related Work}
\label{sec:related}

Our work lies at the intersection of three distinct research areas: multi-provider LLM orchestration, LLM-based evaluation methodology, and systems-level reliability engineering. We survey each in turn, identifying the specific gap that motivates our contribution.

\subsection{Multi-Provider LLM Gateways and Routers}

The rapid proliferation of commercial LLM APIs---each with distinct pricing structures, rate limits, capability profiles, and availability characteristics---has given rise to a class of middleware systems designed to abstract over provider heterogeneity. These multi-provider gateways present a unified API surface to downstream applications while internally managing routing, failover, load balancing, and cost optimization across heterogeneous backends.

\textbf{LiteLLM}~\cite{litellm} provides an open-source OpenAI-compatible proxy layer supporting over 100 LLM providers. It implements automatic failover with configurable retry policies, request-level load balancing, and spend tracking. Its architecture treats each completion request as a stateless transaction: when a provider fails, the identical request is forwarded to an alternative endpoint. LiteLLM does not maintain or reconstruct conversational state across failover boundaries; this is by design, as its abstraction operates at the HTTP request level rather than the dialogue level.

\textbf{Portkey}~\cite{portkey} extends the gateway pattern with a managed control plane offering semantic caching, conditional routing based on request metadata, and detailed observability. Its failover mechanism supports configurable fallback chains with weighted routing strategies. Like LiteLLM, Portkey's failover logic operates per-request: a failed call is retried against an alternative provider with the same payload. The system does not inspect or augment the conversational context during failover. Portkey's documentation explicitly frames its reliability guarantees in terms of uptime and successful response delivery, not conversational coherence.

\textbf{OpenRouter}~\cite{openrouter} provides a commercial routing layer that dynamically selects providers based on cost, latency, and availability. It supports automatic fallback across its provider pool and offers a unified billing interface. OpenRouter's routing decisions are made per-request based on real-time provider health signals. As with the systems above, conversational history is not a first-class object in the routing layer; the system ensures that \textit{a} response is returned, not that the response is informed by the full preceding dialogue.

\textbf{Martian}~\cite{martian} and \textbf{Not Diamond}~\cite{notdiamond} represent a more recent generation of model routers that employ learned routing policies to select the optimal model for a given query based on predicted quality, latency, or cost. These systems introduce sophistication at the model selection layer but inherit the same stateless request-level abstraction: the routing decision is conditioned on the current request, not on the continuity of an ongoing conversation.

Several additional systems merit acknowledgment. \textbf{Amazon Bedrock}~\cite{bedrock} provides managed multi-model inference with cross-region failover capabilities. \textbf{Azure OpenAI Service}~\cite{azure-openai} offers provisioned throughput with automatic regional failover. Both operate within their respective cloud ecosystems and implement failover at the infrastructure level, ensuring request-level availability without addressing dialogue-level continuity.

\textbf{Summary.} Existing multi-provider systems have made substantial and genuine contributions to the operational reliability of LLM deployments. They have largely solved the \textit{availability} problem: ensuring that a well-formed response is returned despite individual provider failures. What they do not address---and what we argue constitutes a distinct, orthogonal property---is \textit{continuity}: ensuring that the response returned after a failover event reflects the full conversational context established with the now-unavailable provider. This paper formalizes that distinction and provides the first systematic evaluation framework for measuring it.

\subsection{LLM Evaluation and LLM-as-Judge}

Our evaluation methodology relies on automated LLM judging to assess whether conversational context has been preserved across failover boundaries. This approach draws on a rapidly maturing body of work on using language models as evaluators.

\textbf{LLM-as-Judge}~\cite{zheng2023judging} demonstrated that strong language models (e.g., GPT-4) can serve as reliable proxies for human evaluation across a range of generation tasks, achieving high agreement with human raters on quality assessments. Subsequent work has refined this paradigm: \textbf{Chatbot Arena}~\cite{chiang2024chatbot} employs pairwise LLM judging at scale to produce Elo-style model rankings from crowd-sourced preferences, while \textbf{AlpacaEval}~\cite{alpacaeval} uses automated LLM comparison against a reference model to approximate human preference judgments.

Our use of LLM-as-judge differs from these efforts in a critical respect. Prior work uses LLM judges to assess open-ended \textit{quality}---fluency, helpfulness, harmlessness---where ground truth is inherently subjective. Our evaluation task is substantially more constrained: the judge must determine whether a specific factual anchor (e.g., a date, a name, a stated preference) established earlier in the conversation is correctly recalled in the post-failover response. This is a \textit{factual verification} task rather than a preference judgment, and we expect---and empirically observe---higher inter-rater reliability as a consequence.

\textbf{FActScore}~\cite{min2023factscore} and \textbf{HaluEval}~\cite{halueval} evaluate factual consistency and hallucination detection in generated text, providing methodological precedent for using LLMs to verify specific factual claims. Our factual anchoring strategy---embedding distinctive, verifiable facts (invented proper nouns, specific dates, unusual preferences) in early conversation turns and probing for their recall after failover---is informed by this line of work but adapted to the novel setting of cross-provider context transfer.

\subsection{Systems Reliability and Failover}

The problem of maintaining service continuity during component failures has deep roots in distributed systems research. \textbf{Circuit breakers}~\cite{nygard2007release}, popularized in microservice architectures, prevent cascading failures by temporarily halting requests to unhealthy downstream services. \textbf{Bulkhead patterns}~\cite{nygard2007release} isolate failure domains to contain blast radius. \textbf{Exponential backoff with jitter}~\cite{aws-backoff} is the standard approach for preventing synchronized retry storms against rate-limited or recovering services, a technique whose necessity we independently rediscover and empirically validate in the LLM failover setting (Section~\ref{sec:discussion}).

The concept of \textbf{session affinity} (or ``sticky sessions'')~\cite{hunt2010zookeeper} in distributed systems---ensuring that sequential requests from the same logical session are routed to the same backend---is a partial analogue to our continuity requirement. However, session affinity in traditional systems ensures routing consistency, not state reconstruction: if the sticky backend fails, the session state is typically lost unless explicitly replicated. Our History-Forwarding strategy can be understood as active state reconstruction at failover time, analogous to \textbf{state machine replication}~\cite{schneider1990implementing} in fault-tolerant distributed systems, but operating on conversational dialogue rather than deterministic state machines.

\textbf{The thundering herd problem}~\cite{bonwick1994thundering} describes the pathology in which a large number of processes are simultaneously awakened to compete for a shared resource, overwhelming it. Our empirically observed retry-storm failure mode (Section~\ref{sec:discussion}) is a direct instantiation of this classical problem in the specific context of LLM API failover under high concurrency, where the ``shared resource'' is a rate-limited fallback provider's request quota.

\subsection{Positioning}

To our knowledge, no prior work has (1) formally defined conversational continuity as a property distinct from availability in multi-provider LLM systems, (2) proposed quantitative metrics for its evaluation, or (3) systematically measured it under controlled failover conditions. This paper addresses all three.
